\subsection{Matriks}
Matriks adalah susunan bilangan-bilangan riil yaang diatur dalam yang diatur dalam baris dan kolom sehingga membentu \textit{array} persegi panjang. Secara formal sebuah matriks $A$ berukuran $m\times n$ berarti suatu kumpulan $m \cdot n$ bilangan yang disusun dalam $m$ baris dan $n$ kolom. 

$$
A = \begin{pmatrix}
a_{11} & a_{12} & \dots & a_{1n} \\
a_{21} & a_{22} & \dots & a_{2n} \\
\vdots & \vdots & \ddots & \vdots \\
a_{m1} & a_{m2} & \dots & a_{mn}
\end{pmatrix}
$$

Matriks dapat ditulis sebagai $A \in \mathbb{R}^{m \times n}$, di mana setiap elemen $a_{ij}$ mewakili bilangan yang terletak di pada baris ke-$i$ dan kolom ke-$j$. Matriks dapat disebut sebagai matriks persegi, ketika baris dan kolomnya memiliki ukuran yang sama ($m = n$). Jika, matriks berbentuk persegi, maka elemen $a_{ij}$ di mana $i = j$ disebut sebagai diagonal utama dari matriks tersebut.

Pada dasarnya, vektor merupakan matriks berukuran $1 \times n$ (vektor baris) atau $n \times 1$ (vektor kolom). Selain itu, terdapat juga bentuk matriks Identitas ($I_n$)dan matriks Nol ($0$). Matriks identitas $I_n$ merupakan matriks persegi $n 
times n$ dengan diagonal utama bernilai 1 dan elemen lainnya 0. Sedangkan, matriks nol adalah matriks dengan semua elemen bernilai 0. Terakhir, terdapat matriks diagonal, yaitu matriks persegi dengan elemen non-nol hanya pada diagonal utama. 

$$
I_3 = 
\begin{pmatrix}
1 & 0 & 0 \\
0 & 1 & 0 \\
0 & 0 & 1 \\
\end{pmatrix}
\hspace{0.3cm}
0 = \begin{pmatrix}
    0   & 0     & 0 \\
    0   & 0     & 0 \\
    0   & 0     & 0 \\
\end{pmatrix}
\hspace{0.3cm}
D = \begin{pmatrix}
    d_1 & 0     & 0 \\
    0   & d_2   & 0 \\
    0   & 0     & d_3 \\
\end{pmatrix}
$$

\subsection{Transpose Matrix}
Transpose dari sebuah matriks $A \in \mathbb{R}^{m \times n}$ dinotasikan sebagai $A^T$ adalah matriks berukuran $n \times m$ yang diperoleh dengan menukar baris dan kolom dari A. 
\begin{align*}
    A = (a_{ij}) &\implies A^T = (a_{ji}) \\
    A = \begin{pmatrix}
        a_{11} & a_{12}  & a_{13} \\
        a_{21} & a_{22}  & a_{23}
    \end{pmatrix}_{2 \times 3}
                 &\implies
    A^T = \begin{pmatrix}
        a_{11} & a_{21} \\
        a_{12} & a_{22} \\
        a_{13} & a_{23} 
    \end{pmatrix}_{3 \times 2}
\end{align*}

\subsection{Matriks Simetris}
Sebuah matriks persegi A disebut simetris jika $A = A^T$ artinya matriks sama dengan transposenya ($a_{ij} = a_{ji}$).
Setiap matriks $AA^T$ atau $A^TA$ yang perkaliannya terdefinisi merupakan matriks simetris.
Sebuah matriks simstris memiliki properti sebagai berikut:
\begin{IEEEitemize}
    \item Semua nilai eigen dari matriks simetris adalah bilangan riil
    \item Vektor eigen yang bersesuaian dengan nilai eigen yang berbeda adalah orthogonal
    \item Matriks simetris $A$ selalu dapat didiagonalisasi sehingga $A = Q \varLambda Q^T$ di mana $Q$ adalah matriks orthogonal dan $\varLambda$ adalah matriks diagonal.
\end{IEEEitemize}

\subsection{Perkalian Matriks}
Perkalian dua matriks $A \in \mathbb{R}^{m \times n}$ dan $B \in \mathbb{R}^{n \times p}$ menghasilkan matriks $C \in \mathbb{R}^{m \times p}$ dengan elemen-elemen:
$$
C_{ij} = \sum{k = 1}^n A_ik \times B_kj
$$
Jumlah kolom $A$ harus sama dengan jumlah baris $B$ agar perkalian terdefinisi.

\subsection{Hadamard Product}
Perkalian \textit{element-wise} antara dua matriks $A$ dan $B$ dengan ukuran sama $m \times n$ menghasilkan matriks $C$ dengan elemen-elemen
$$
C = A  \odot B \implies C_{ij} = A_{ij} \cdot B_{ij}
$$ 

\subsection{Trace}
Trace dari matriks persegi $A \in \mathbb{R}^{n \times n}$ adalah penjumlahan elemen-elemen diagonal utamanya.
$$
\mathrm{tr}(A) = \sum{i = 1}^n a_{ii} = a_{11} + a_{22} + \dots + a_{nn}
$$

\subsection{Frobenius Norm}
Frobenius norm digunakan untuk mengukur \textit{magnitude} atau ukuran dari matriks secara keseluruhan. Ini merupakan generalisasi dari Euclidean norm untuk vektor. 
Frabernius norm dari matriks $A \in \mathbb{R}^{m \times n}$ dapat dihitung dengan:
$$
\Vert A \Vert_F = \sqrt{\sum{i = 1}^m \sum{j = 1}^n a{ij}^2} = \sqrt{\mathrm{tr}(A^TA)}
$$
